% !TeX encoding = UTF-8
% !TeX root = thesis.tex

\section{Anhang}

Die zur Reproduktion relevanten Artefakte befinden sich im Repository: die vier
Jupyter-Notebooks (\texttt{01\_data\_exploration}, \texttt{02\_metric\_baseline},
\texttt{03\_log\_component}, \texttt{04\_fusion}), die erzeugten Abbildungen
im Verzeichnis \texttt{docs/} sowie die 13~Unit-Tests in
\texttt{tests/test\_pipeline.py}.

Ergänzend liegt im Repository ein \texttt{Dockerfile} vor, das eine einfache
Jupyter-Lab-Umgebung auf Basis von Python~3.12 mit den benötigten Abhängigkeiten
bereitstellt. Es dient der lokalen Reproduzierbarkeit der Notebook-Ausführung und
ist nicht als produktionsreife Deployment-Umgebung zu verstehen. Der Container kann
aus dem Repository-Wurzelverzeichnis mit
\texttt{docker build -t cloudanobench-thesis .}
gebaut und anschließend mit
\texttt{docker run -p 8888:8888 cloudanobench-thesis}
gestartet werden.
